\section{VC 维把无限假设类变有限}
\label{sec:13-Machine-Learning/10-Learning-Theory/03}

你手头有一个假设类 \(\mathcal{H}\)，里面装着无限多个函数。上一节的联合界给了你一个泛化误差上界，里面有一项 \(\log|\mathcal{H}|\)。有限类的时候你心安理得地把它代进去——三千个函数，\(\log 3000\approx 8\)，完全可以接受。现在类里的函数个数是无穷大，\(\log\infty\) 没法用。你第一反应可能是：数一下参数个数？我的模型有 50 万个参数，用参数个数替代 \(|H|\) 行不行？

下面三分钟之内你会看到为什么这个直觉是死路，以及一条完全不同的出路。

\subsection{为什么数参数不解决问题}

先试一个看起来温和的类：实轴上的区间分类器，
\[
\mathcal{H}=\{\ind\{x\in[a,b]\}\mid a\le b\}.
\]
两个参数，\(a\) 和 \(b\)。在 \(n\) 个点上它最多能产生多少种不同的标记方式？想象点按大小排好，一个区间无非是取一个连续段标 \(+1\)、其余标 \(-1\)。在 \(n\) 个点中选连续段，起止位置各有 \(n+1\) 种选择（包括``不取任何点''），所以不同标记方式大约是 \(O(n^2)\)——比 \(2^n\) 少得多。参数个数是 2，这碰巧对得上某种直觉。

再看 \(\mathbb{R}^d\) 中的线性分类器，\(d+1\) 个参数，VC 维也是 \(d+1\)。参数个数 = VC 维，似乎一条规律正在浮现。

现在看第三个类：
\[
\mathcal{H}=\{x\mapsto \sign(\sin(\omega x))\mid \omega>0\}.
\]
只有一个参数 \(\omega\)。直觉告诉你：一个参数只能拧出一个自由度，行为应该很受限。我们来找 \(n\) 个点。取 \(x_i=10^{-i}, i=1,\dots,n\)。这串点以十倍的速度缩向 0，相邻点之间的距离每步缩小一个数量级。任给一组标记 \(y_1,\dots,y_n\in\{-1,+1\}\)，我需要找到一个 \(\omega\) 使得 \(\sign(\sin(\omega x_i))=y_i\) 对每个 \(i\) 成立。

构造方法不复杂但需要一点数论直觉。令
\[
\omega = \pi\left(\sum_{i=1}^n \frac{1-y_i}{2}\,10^i + \frac12\right).
\]
考察 \(\omega x_j/\pi\) 除以 2 后的余数。\(i>j\) 的项被 \(10^{i-j}\) 放大为偶整数，在模 2 意义下消失；\(i=j\) 的项贡献 \((1-y_j)/2\)，即 0 或 1；\(i<j\) 的项总和不超过 \(\sum_{k\ge 1}10^{-k}=1/9\)，远不够改变余数所属的整数区间；末尾的 \(1/2\) 再加一个小的正偏移。结果是：\(y_j=+1\) 时 \(\omega x_j/\pi\) 的余数严格在 \((0,1)\) 内，正弦为正；\(y_j=-1\) 时余数严格在 \((1,2)\) 内，正弦为负。

这套操作对任意大的 \(n\) 都成立。一个参数的类打散了任意多个点。VC 维 \(=+\infty\)。

三个类排在一起，教训是清楚的：VC 维度量的是类在有限样本上的组合表现力，不是参数个数，更不是集合基数。

\subsection{换一个问题：在 \(n\) 个点上能折腾出多少花样}

既然``类里有多少函数''问了错的问题，正确的问题是什么？回到泛化界的证明逻辑：联合界之所以需要 \(\log|\mathcal{H}|\)，是因为我们要同时约束类中所有函数的泛化误差，而有限类的做法是对每个函数单独用 Hoeffding 再取并集。真正起作用的不该是宇宙中这个类有多大，而应该是——给定当前这 \(n\) 个数据点，类中的函数在这些点上能表现出多少种不同的行为。

考虑二分类 \(h:\mathcal{X}\to\{-1,+1\}\)。给定 \(n\) 个点 \(x_1,\dots,x_n\)，每个假设在这批点上留下的只是一个标记向量
\[
\bigl(h(x_1),\dots,h(x_n)\bigr)\in\{-1,+1\}^n.
\]
两个在这批点上标记完全相同的假设，对这批数据而言无法区分——你无论用哪个，训练误差和泛化行为在这一批数据的意义上没有区别。所以``类的有效大小''该用去重之后的等价类个数来衡量。

这引出生长函数（growth function）：
\[
\Pi_{\mathcal{H}}(n)
= \max_{x_1,\dots,x_n}
\abs{\bigl\{\bigl(h(x_1),\dots,h(x_n)\bigr)\mid h\in\mathcal{H}\bigr\}}.
\]
取 max 意味着我们考虑的是最坏情形——在所有可能的 \(n\) 点配置中，选那个让类能表现出最多不同行为的那组点。这是悲观但安全的做法：泛化界要对任意数据分布成立，自然要对最坏的点配置也成立。

显然 \(\Pi_{\mathcal{H}}(n)\le 2^n\)，因为 \(n\) 个点的二分类标记总共就 \(2^n\) 种。如果等号成立——这批点的每一种标记都能被类中某个假设实现——我们称 \(\mathcal{H}\) 打散（shatter）了这批点。

打散是本节的枢纽概念。它不关心类里有多少函数，只关心类在具体一批点上的标记能力是否``满格''。

\subsection{VC 维：打散能力的临界规模}

一个类也许能打散 3 个点、5 个点，但终归会遇到一个上限——再多的点就无论如何排列，也不可能打散。这个临界规模就是 VC 维：

\[
d_{\mathrm{VC}}(\mathcal{H})
= \max\{n\mid \Pi_{\mathcal{H}}(n)=2^n\}.
\]

定义里的 max 意味着：只要存在某组 \(n\) 个点被打散，我们就说 VC 维至少是 \(n\)。不需要所有 \(n\) 点集合都被打散。VC 维是类的最强标记能力声明——``在最有利的点配置下，它能完全自由地标到多大规模''。

回到那三个例子验证一下。区间分类器：任取两点 \(x_1<x_2\)，四种标记全能实现（区间只包 \(x_1\)、只包 \(x_2\)、两个都包、两个都不包），所以 \(d_{\mathrm{VC}}\ge 2\)。任取三点 \(x_1<x_2<x_3\)，标记 \((+,-,+)\) 不可能是区间——含两个端点的区间必然包含中点。所以 \(d_{\mathrm{VC}}=2\)。\(\mathbb{R}^d\) 中线性分类器：处于一般位置的 \(d+1\) 个点总可以被半空间实现任意标记，下界 \(d+1\)；Radon 定理保证任意 \(d+2\) 个点可分成两组使凸包相交，把两组标以相反符号则任何半空间都无法分离，上界 \(d+1\)。正弦类：我们刚构造了能打散任意 \(n\) 个点的方法，\(d_{\mathrm{VC}}=+\infty\)。

\subsection{Sauer 引理：一旦不能打散，生长立刻从指数塌成多项式}

生长函数最引人注目的性质是它只有两种命运。要么对所有 \(n\) 都有 \(\Pi_{\mathcal{H}}(n)=2^n\)——这意味着 \(d_{\mathrm{VC}}=+\infty\)，类能打散任意大的点集。要么一旦 \(n\) 超过 \(d=d_{\mathrm{VC}}(\mathcal{H})\)，生长立刻减速。

减速有多剧烈？Sauer 引理给出确切的答案：
\[
\Pi_{\mathcal{H}}(n)
\le \sum_{i=0}^{d}\binom{n}{i}
\le \left(\frac{en}{d}\right)^{d},
\qquad n\ge d.
\]
第一个不等式说：能实现的标记方式不会多于``从 \(n\) 个点中选不超过 \(d\) 个子集的组合方式''。第二个不等式是宽松但更好记的形式。取对数：
\[
\log\Pi_{\mathcal{H}}(n) \le d\log\frac{en}{d} \approx d\log n.
\]

这就是``无限假设类在有限数据上其实很小''的精确含义。类中的函数或许在集合论意义上不可数，但它们在任意 \(n\) 个训练点上能表现出的不同行为只有 \(n\) 的多项式种——而非指数种。一旦 \(n\) 超过 VC 维，类就不再有足够的自由度去遍历所有可能的标记组合，只有一些``特殊''的标记能被实现，而这些特殊标记的数量被多项式控制住了。

你可能会想：一个不可数集怎么会在有限点上只有多项式种行为？因为行为空间的大小被数据点数锁死了——\(n\) 个点最多容纳 \(2^n\) 种标记，而一旦 VC 维有限，这个 \(2^n\) 的上界就永远达不到，实际能达到的标记数被压进了多形式。

\subsection{VC 界：把 \(\log|\mathcal{H}|\) 换成生长函数的对数}

现在把刚才的直觉接回到泛化界。联合界中使用 \(\log|\mathcal{H}|\) 的位置，用生长函数的对数替换（证明的技术细节——对称化引入幽灵样本、用 \(2n\) 个点枚举行为——在此略去），就得到 VC 泛化界：以至少 \(1-\delta\) 的概率，对所有 \(h\in\mathcal{H}\) 同时，
\[
R(h)\le\widehat R(h)
+ O\left(
\sqrt{\frac{d_{\mathrm{VC}}\log\frac{n}{d_{\mathrm{VC}}}+\log\frac{1}{\delta}}{n}}
\right),
\]
其中 \(d_{\mathrm{VC}}=d_{\mathrm{VC}}(\mathcal{H})\)。结构上与有限类的界完全一致，只是``有效大小的对数''从 \(\log|\mathcal{H}|\) 变成了 \(d_{\mathrm{VC}}\log n\)。

这个界直接回答了``需要多少样本''。沿用 ERM 论证——超出风险不超过两倍界宽——要以至少 \(1-\delta\) 的概率把超出风险压到 \(\varepsilon\) 以内，样本量需要
\[
n \sim \frac{d_{\mathrm{VC}}\log(1/\varepsilon) + \log(1/\delta)}{\varepsilon^2}
\]
量级。样本量随 VC 维近似线性增长，与 \(|\mathcal{H}|\) 无关。

现在我们可以回答本篇的核心问题了：什么样的假设类在分布无关的意义下可学？VC 维有限的类——无论基数多大——可学。VC 维无限的类——哪怕只有一个参数——在分布无关的意义下不可学：没有任何有限样本量能保证在所有分布上同时把泛化误差压到 \(\varepsilon\) 以内。正弦类的教训是，哪怕一个参数，如果它能制造无限复杂的决策边界，有限数据就永远追不上它。

这不是凭空的哲学判断，而是可以从 VC 界严格写出的结论。

\subsection{证明在做什么，以及为什么 Rademacher 复杂度更好}

这里不展开 VC 界的完整证明，但值得记住它的两步骨架，因为第二步本身就演化成了比生长函数更紧的现代工具。

第一步，对称化（symmetrization）。引入一份``幽灵样本''（ghost sample）——从同一个分布独立抽取的另 \(n\) 个点。把真风险 \(R(h)\) 写成在这份幽灵样本上的经验风险期望，于是 \(R(h)-\widehat R(h)\) 变成两份经验风险的差。分布的难题由此转化为两份有限样本的比较问题，可以在 \(2n\) 个点上枚举行为。生长函数正是在这里登场的——它算出在这 \(2n\) 个点上最多有多少种行为模式需要被 union bound 覆盖。

第二步，用随机符号替代实际标记差。给每个样本注入一个独立的 Rademacher 随机变量（以等概率取 \(\pm 1\)），把``类能实现多少种标记''归约为``类拟合纯随机噪声的能力''。这一步的巧妙在于：用随机性消除了对真实标记分布的依赖，留下的量只取决于数据点在输入空间中的相对几何位置。这就是 Rademacher 复杂度。

Rademacher 复杂度比生长函数更紧，因为它把最坏情形的组合计数换成了数据依赖的平均拟合噪声能力。生长函数是对所有 \(n\) 点配置取 max——它永远按最坏的布局来算；Rademacher 复杂度只对当前手里的这批数据取期望。前者是分布无关的、纯组合的；后者是数据相关的、带概率的。下一篇范数界正是从 Rademacher 复杂度出发，用权重范数替代 VC 维作为复杂度度量。

这里只需记住它在论证链中的位置：VC 维和生长函数完成了从``无穷假设类无法用联合界''到``多项式样本复杂度''的关键跨越，而 Rademacher 复杂度在同一框架内提供了更精细的刻度。
